%==============================================================================
%  CONFERENCE PAPER (D49): the fixed-threshold transfer audit. fp32, desktop.
%  Plan: OneDrive phase_4_writing/step_05_companion_paper.md. Split rationale:
%  02_protocol/publication_split.md. Do NOT add conversion, int8, QTS (H3), H6,
%  ExecuTorch or device results: they are the journal paper's new content.
%  Venue-neutral IEEE two-column draft, at most 6 pages including references.
%  Every number carries a "% src:" comment naming the committed file it comes
%  from (repo-relative). Prose is written fresh: the TIFS draft
%  (paper/tifs_manuscript.tex) shares numbers and figures, not paragraphs.
%  Build: pdflatex main; bibtex main; pdflatex main; pdflatex main
%==============================================================================
\documentclass[conference]{IEEEtran}
\usepackage{cite}
\usepackage{amsmath,amssymb}
\usepackage{graphicx}
\graphicspath{{figures/}}
\usepackage{booktabs}
\usepackage{xcolor}
\usepackage{url}

% ---- review mode ------------------------------------------------------------
\newif\ifblind
\blindfalse   % \blindtrue for double-blind review: anonymous author block and links
\newcommand{\todo}[1]{\textcolor{red}{[#1]}}
\ifblind
  \newcommand{\repo}{\todo{anonymised repository link}}
  \newcommand{\preregref}{\todo{anonymised link}}
\else
  \newcommand{\repo}{\todo{public repository URL}}
  \newcommand{\preregref}{tag \texttt{roster-freeze}, commit \texttt{97e9753}}
\fi

% ---- notation ---------------------------------------------------------------
\newcommand{\ttg}{\mathrm{TTG}}
\newcommand{\hter}{\mathrm{HTER}}
\newcommand{\frr}{\mathrm{FRR}}
\newcommand{\far}{\mathrm{FAR}}
\newcommand{\eer}{\mathrm{EER}}
\newcommand{\tauE}{\tau^{*}_{\mathrm{EER}}}
\newcommand{\tauF}[1]{\tau^{*}_{\mathrm{FRR}#1}}
\newcommand{\tauor}{\tau_{\mathrm{or}}}

\begin{document}

% Working title. The title is a [U] decision (D30); alternatives in progress.md.
\title{Calibrated Once, Deployed Elsewhere:\\ A Pre-Registered Audit of
Threshold Transfer in Speech Deepfake Detection}

\ifblind
\author{\IEEEauthorblockN{Anonymous submission}}
\else
\author{\IEEEauthorblockN{Hoan Hoang}
\IEEEauthorblockA{\todo{affiliation}\\
\todo{e-mail}}}
\fi

\maketitle

% ============================ ABSTRACT =======================================
% src (every number): results/analysis/operating_points.txt (FRR at tau*_FRR1 on ITW:
% 42.78-81.78; screening C > 1 for 5/5 models on ITW and DF21 at all policies; XLS-R ITW
% FRR 49.90); h2_test.txt (3/4); h5_test.txt (4/4, FAR falls for all four); h1_h4.txt
% (rho 0.00 [-0.62, 0.58]); paper/figures/figures_data.json (XLS-R ITW safe band 0.034 SD).
\begin{abstract}
Speech deepfake countermeasures are compared by equal error rate (EER), a
threshold-free measure, but a deployed detector must fix its threshold in
advance. We measure what that commitment costs for five published
countermeasures: four ranked models (AASIST, AASIST-L, RawNet2, Res-TSSDNet)
and an XLS-R self-supervised reference. Thresholds frozen on ASVspoof~2019 LA
development data at three operating points are transferred unchanged to the
2019 LA evaluation set, In-the-Wild, ASVspoof~2021 DF and neural-codec
resynthesis, under hypotheses registered before any In-the-Wild score existed.
We report the threshold transfer gap with its error decomposition and a
threshold-sensitivity safe band. A 1\% false-rejection budget set on source
data becomes 43--82\% on In-the-Wild. There, false rejection exceeds false
acceptance for three of four ranked models, and EnCodec resynthesis raises
false rejection for all four while false acceptance falls (both pre-registered
tests). Source EER shows no detectable association with the transfer gap
(\(\rho=0.00\), 95\% CI \([-0.62,0.58]\)). At its transferred threshold, the
most accurate model rejects half of all genuine In-the-Wild speech; its
saturated scores leave a safe band only 0.03 standard deviations wide. Under a
screening cost model, every transferred detector costs more than no detector on
both shifted corpora. We release the harness.
\end{abstract}

\begin{IEEEkeywords}
speech deepfake detection, anti-spoofing, decision threshold, calibration,
domain shift, neural audio codec, pre-registration
\end{IEEEkeywords}

% ============================ 1. INTRODUCTION ================================
\section{Introduction}
\label{sec:intro}
Speech deepfake countermeasures (CMs) are ranked by the equal error rate (EER)
and the minimum tandem detection cost \cite{kinnunen18_tdcf}, both evaluated at
the threshold that suits the evaluation data best
\cite{wang20_asvspoof2019,liu23_asvspoof2021}. Both summarise how well a CM
ranks a test set. A deployed CM needs more than a ranking. A call-screening
filter, a moderation queue or a browser plug-in \cite{pascu25_edge_plugin}
commits to a threshold before deployment, chooses it on data it owns, and then
makes hard decisions on traffic whose distribution it does not control. The EER
says nothing about how a threshold chosen elsewhere will behave there, and the
difference can be large. Zhou and Wang \cite{zhou26_eer_hides} report a detector with
0.21\% source EER that rejects 78.7\% of genuine In-the-Wild speech once its
source threshold is transferred.
% src: zhou26 (arXiv:2606.21584v1) Table 1, verified against the PDF in OneDrive
% 01_references/notes/tier1/zhou26_eer_hides.md (tau set on the LA19 EVAL set, one model).

That measurement covers one model at one operating point, with the threshold
set on the labelled source evaluation set. A deployer needs three further
answers. Is the failure general across architectures and targets, and does it
have a direction? Does a lower EER protect a detector? Does a transmission
channel move the operating point in the same way? We address these questions
for five published CMs: four ranked models (AASIST and AASIST-L
\cite{jung22_aasist}, RawNet2 \cite{tak21_rawnet2}, Res-TSSDNet
\cite{hua21_res_tssdnet}) and a self-supervised (SSL) reference, XLS-R+AASIST
\cite{tak22_xlsr_aasist}. Checkpoints are used as released. Thresholds are
frozen on ASVspoof~2019 LA (LA19) development data and transferred unchanged to
the LA19 evaluation set, In-the-Wild (ITW) \cite{muller22_in_the_wild}, a
verified-audio surface of ASVspoof~2021 DF (DF21) \cite{liu23_asvspoof2021} and
EnCodec-resynthesised speech \cite{defossez23_encodec}. The hypotheses were
registered, and frozen, before any ITW score existed. Our contributions:
\begin{itemize}
\item \emph{A fixed-threshold transfer protocol}: the threshold transfer gap
  (TTG) with a minimum-HTER reference, defined at any operating point; a
  mandatory false-rejection/false-acceptance (FRR/FAR) decomposition; a
  threshold-sensitivity curve with a 1-pp safe band; and a registered
  cost-weighted robustness check. The harness is released.
\item \emph{A pre-registered five-model audit} with two confirmatory tests under
  Holm correction (H2: FRR exceeds FAR on ITW; H5: neural-codec resynthesis
  raises FRR more than FAR) and two exploratory estimates (H1: source EER
  against TTG; H4: model size against TTG).
\item \emph{Four findings.} Unseen attacks raise false acceptance, but under
  corpus shift the larger error is genuine speech rejected, and a
  false-rejection budget fixed on source data does not survive transfer.
  Source EER shows no detectable association with the gap, and the most
  accurate model has the largest one, which its saturated scores explain. A
  neural codec lowers the scores of genuine and fake speech alike. Which model
  loses most depends on the deployment cost model.
\end{itemize}
All scores here are fp32 on a desktop. Conversion, quantization and on-device
behaviour of the same models are reported separately.

% ============================ 2. RELATED WORK ================================
\section{Related Work}
\label{sec:related}
\emph{Evaluation practice.} ASVspoof ranks CMs by EER and minimum t-DCF
\cite{wang20_asvspoof2019,liu23_asvspoof2021,kinnunen18_tdcf}. ASVspoof~5 also
reports actual costs at a fixed threshold, and its strongest systems pair low
minimum costs with actual costs near 1 \cite{wang24_asvspoof5}, a calibration
failure that EER cannot show. Speaker recognition separates discrimination from
calibration \cite{brummer06_calibration}, and biometric authentication reports
the half total error rate (HTER) at a threshold fixed in advance
\cite{bengio04_hter}. We carry that practice over to CMs.

\emph{Generalization.} ITW \cite{muller22_in_the_wild} showed that CMs trained
on ASVspoof data degrade sharply on genuine and fake speech of public figures
collected from online media. MLAAD \cite{mlaad24} adds a language axis, and
Berisha et al.\ \cite{berisha25_coverage_debt} argue that deployment conditions
multiply faster than training data can cover them. These studies report the
target EER, that is, an oracle operating point.

\emph{Threshold transfer.} The closest precedent is Zhou and Wang
\cite{zhou26_eer_hides}. They freeze one SSL detector and transfer its
source-EER threshold to ITW and a DF21 subset. They show that monotone,
unlabelled score corrections cannot change the ranking, and that the one family
that moves the operating point helps on one target and fails on another that
differs in class prior. Their threshold comes from the labelled source
evaluation set. Ours come from development data, so no evaluation label
informs any threshold. We add five models, three operating points, the full ITW
corpus, a verified DF21 surface, pre-registered tests, and a sensitivity
analysis that explains why the gaps differ.

\emph{Channels.} DF21 builds media compression into its evaluation
\cite{liu23_asvspoof2021}, and communication benchmarks and challenges add
VoIP and media transformations \cite{realworld_comms24,radar26}. Neural codecs
are also a synthesis route: CMs trained on vocoder-era data miss codec-based
fakes \cite{wu24_codecfake}. These studies score at oracle thresholds. We ask
what a neural codec does to a fixed operating point, for genuine speech as well
as for fakes.

% ============================ 3. PROTOCOL =====================================
\section{Protocol}
\label{sec:protocol}

\subsection{Score contract and frozen thresholds}
% src: results/phase_1_reproduction/tau_<m>.json (dev n = 24,844; frozen_utc fields);
% dev bona fide count 2,548 (repro_report_<m>.md). ITW first scored 2026-09-29 and DF21
% for every model after its tau was frozen (progress.md).
Each CM maps a clip \(x\) to a score \(s(x)\), higher meaning more bona fide,
and declares it a spoof if \(s(x)<\tau\). All five models were trained on the
LA19 training set \cite{wang20_asvspoof2019}, so every threshold is fixed on the
LA19 development set (24,844 trials, 2,548 bona fide) at three operating
points: \(\tauE\), where dev FRR equals FAR, and \(\tauF{1}\) and \(\tauF{5}\),
the strictest thresholds whose dev FRR stays within 1\% and 5\%. The two FRR
policies encode a deployer's budget for rejecting genuine users. Each threshold
is computed from dev scores alone and frozen at reproduction time, before any
ITW or DF21 score existed. It is then applied unchanged to every target.
Nothing is re-calibrated.

\subsection{Measures}
On a target, \(\frr(\tau)\) is the share of bona fide clips declared spoof,
\(\far(\tau)\) the share of spoofs accepted, and
\(\hter(\tau)=\tfrac12[\frr(\tau)+\far(\tau)]\) \cite{bengio04_hter}. For a
transferred threshold \(\tau^{*}\), the threshold transfer gap
\begin{equation}
\ttg(\tau^{*}) \;=\; \hter(\tau^{*}) \;-\; \min_{\tau}\hter(\tau)
\label{eq:ttg}
\end{equation}
is the error added by using \(\tau^{*}\) instead of the best fixed threshold
for that target. Taking the minimum HTER rather than the EER as reference makes
TTG well defined at every operating point. At the EER point, the analogous
quantity can be read from \cite{zhou26_eer_hides} for one model. We always
report TTG with its FRR/FAR decomposition, because equal gaps can hide opposite
failures.

\emph{Threshold sensitivity.} We trace \(\hter(\tau)\) on each target and
express \(\tau\) in standard deviations (SD) of the target scores, measured from
the target-oracle threshold \(\tauor=\arg\min_\tau \hter(\tau)\). This puts models
with different score scales on one axis. Two numbers summarise each curve. The
\emph{miss} is \((\tau^{*}-\tauor)/\mathrm{SD}\). A positive miss means the
transferred threshold is stricter than the target optimum, so its excess error
is mostly false rejection. The \emph{safe band} is the width, in SD, of the
threshold range in which HTER stays within 1\,pp of its minimum. A narrow band
makes a small miss expensive.

\emph{Cost-weighted check (registered).} HTER weights both errors equally. The
registration adds two deployment models: \emph{screening}, with bona fide prior
\(\pi=0.95\) and unit costs, and \emph{forensic}, with \(\pi=0.5\) and a false
acceptance ten times as costly as a false rejection (\(c_{\mathrm{FA}}=10\),
\(c_{\mathrm{FR}}=1\)). For each we report the normalised cost
\begin{equation}
C(\tau)=\frac{\pi\,c_{\mathrm{FR}}\,\frr(\tau)+(1-\pi)\,c_{\mathrm{FA}}\,\far(\tau)}
{\min\left[\pi\,c_{\mathrm{FR}},\;(1-\pi)\,c_{\mathrm{FA}}\right]}.
\label{eq:cost}
\end{equation}
The denominator is the cost of the better trivial system (accept everything or
reject everything), so \(C>1\) means that the detector costs more than not
detecting at all.

\subsection{Models}
\label{sec:models}
% src: results/analysis/param_counts.json (params); results/analysis/h1_h4.txt notes (XLS-R
% 315.4M inference graph); models/base.py (input lengths); results/phase_1_reproduction/
% repro_report_<m>.md (EERs, anchors, gates); repro_report_rawnet2_df21clean.md (15.0636% vs
% 15.0640%); decision.md D33 (per-trial max |d| 3.7e-5 over 305,525 trials); D31 (XLS-R ckpt).
Table~\ref{tab:roster} lists the roster. We vendor each author's code and
reproduce its scoring: the authors' score function and class order, and 16\,kHz
input of 64,600 samples (96,000 for Res-TSSDNet), with short clips repeat-padded
and long ones truncated from the start. Each ranked model must reproduce its
published LA19 evaluation EER to within 15\% (relative). The released RawNet2
checkpoint is the DF21 baseline, a different training run from its published
LA19 figure, so it is gated on DF21 instead: on the 305,525 cleanly decodable
trials (Sec.~\ref{sec:targets}) its EER equals the organisers' 15.06\%, and
per-trial scores agree to within \(3.7{\times}10^{-5}\). The XLS-R model is the
authors' DF-track checkpoint, trained with RawBoost augmentation
\cite{tak22_rawboost}. No LA19 evaluation figure is published for it, so it
enters ungated, as an unranked reference outside every hypothesis count.

\begin{table}[!t]
\caption{Roster: four ranked countermeasures and one SSL reference, all
trained on LA19 train. EER on LA19 eval (\%): ours and published.}
\label{tab:roster}
\centering
\footnotesize
\setlength{\tabcolsep}{4pt}
\begin{tabular}{@{}llrrr@{}}
\toprule
Model & Architecture & Params & Ours & Publ. \\
\midrule
AASIST \cite{jung22_aasist}            & SincConv + graph attn. & 298K  & 0.83 & 0.83 \\
AASIST-L \cite{jung22_aasist}          & SincConv + graph attn. & 85K   & 0.99 & 0.99 \\
RawNet2 \cite{tak21_rawnet2}           & SincConv + GRU         & 17.6M & 4.60 & 5.64\textsuperscript{a} \\
Res-TSSDNet \cite{hua21_res_tssdnet}   & 1-D ResNet             & 349K  & 1.62 & 1.66 \\
XLS-R+AASIST \cite{tak22_xlsr_aasist}  & SSL + graph attn.      & 315M  & 0.15 & ---\textsuperscript{b} \\
\bottomrule
\end{tabular}\\[2pt]
\parbox{\columnwidth}{\footnotesize\textsuperscript{a}Different training run;
gated on DF21 instead (Sec.~\ref{sec:models}). \textsuperscript{b}None published;
ungated and unranked.}
% src: EERs 0.8294/0.9909/4.5955/1.6179 and anchors 0.83/0.99/5.64/1.659 from
% results/phase_1_reproduction/repro_report_<m>.md; XLS-R 0.15 from results/analysis/h1_h4.txt;
% params from results/analysis/param_counts.json (XLS-R: 315.4M inference graph, h1_h4.txt notes).
\end{table}

\subsection{Targets}
\label{sec:targets}
% src: results/analysis/operating_points.json (n, n_bona per target); 58 public figures:
% muller22 (notes/tier1/muller22_in_the_wild.md); nine DF21 conditions, C1 = none: liu23 /
% notes/tier1/yamagishi21_asvspoof2021.md; actDCF ~ 1 vs minDCF 0.075-0.24 (Sec. II):
% notes/tier1/wang24_asvspoof5.md; PREREGISTRATION.md D37
% amendment (A07-A19; A16/A19 = A04/A06); OneDrive phase_1_reproduction/df21_flac_integrity.md
% (228,403 of 533,928 rejected by libFLAC 1.4.3; 78,420 MD5 mismatch); decision.md D34;
% DEVIATIONS.md DEV-2 (H5 source set, EnCodec settings); results/analysis/h5_test.txt (<= 9.5e-7).
\emph{LA19 eval} (71,237 trials, 7,355 bona fide) comes from the same source
corpus as the dev set, with disjoint speakers, but contains attacks absent from
training and development (A07--A19; A16 and A19 reuse the algorithms of A04 and
A06). It is a seen corpus with unseen attacks. \emph{ITW}
\cite{muller22_in_the_wild} holds 31,779 clips (19,963 bona fide) of 58 public
figures, so its genuine speech comes from recording conditions that LA19 never
saw. \emph{DF21} \cite{liu23_asvspoof2021} applies nine media-compression
conditions (one of them uncompressed) to bona fide and spoofed speech. Its
public release carries a decoding hazard. The current reference FLAC decoder (libFLAC~1.4) rejects 43\% of the
533,928 evaluation files. A 2021-era decoder reads them, but for 78,420 files
the decoded audio fails the file's own STREAMINFO MD5, so the audio scored in
2021 cannot be recovered from the public files. We exclude those files and
score a verified surface of 455,508 trials (12,067 bona fide), checking every
fallback decode against its MD5. The 305,525 trials that decode without issue
form a sensitivity surface. Published DF21 figures are therefore not comparable
with ours. \emph{Codec stress (H5).} A frozen 1,000-clip LA19 eval manifest
(500 bona fide, 500 spoof) is resynthesised by EnCodec's 24\,kHz model at
6\,kbps \cite{defossez23_encodec} and resampled to 16\,kHz. A clean copy passes
through the same path, and its scores match the reference scores to within
\(9.5{\times}10^{-7}\). All scores are fp32 PyTorch, with TF32 disabled on GPU.
Every score file, frozen threshold and analysis script is versioned with the
harness, and each analysis refuses to run if a recomputed number disagrees with
a committed one.
% src (TF32): models/base.py load_model (allow_tf32 = False), in place before any GPU score.
% src (refusal): analysis/operating_points.py, analysis/figures/conference_figures.py.

\subsection{Hypotheses, statistics and pre-registration}
% src: PREREGISTRATION.md (H1, H2, H4, H5 statements; frozen decisions); DEVIATIONS.md DEV-2;
% analysis/h2_test.py, h5_test.py, h1_h4.py (B = 1000, seed 0, stratification, Holm).
\emph{H2} (confirmatory): on ITW at \(\tauE\), FRR exceeds FAR for a majority
(at least three) of the four ranked models. \emph{H5} (confirmatory): EnCodec
resynthesis raises FRR more than FAR, relative to the clean copy, for a majority
of the ranked models. \emph{H1} (exploratory): the Spearman correlation between
LA19 evaluation EER and ITW TTG over the four ranked models and three policies.
\emph{H4} (exploratory): the correlation between \(\log_{10}\) parameter count
and TTG at \(\tauE\), per target. The confirmatory tests are one-sided,
class-stratified paired bootstraps \cite{koehn04_bootstrap} (\(B=1000\), seed 0)
of \(\frr-\far\) (H2) and \(\Delta\frr-\Delta\far\) (H5), with Holm correction
\cite{holm79_holm} within each family at \(\alpha=0.05\). TTG intervals come from a
class-stratified bootstrap, and correlation intervals from a bootstrap over
points. Exploratory results carry effect sizes and intervals only.

\emph{Pre-registration statement.} The hypotheses were registered on 19~July
2026, amended in August and September, and frozen on 28~September 2026
(\preregref). Every ITW score and all EnCodec audio postdate the freeze. Of the
DF21 scores, only RawNet2's, computed for its reproduction gate, predate it.
This paper reports H1, H2, H4 and H5. The other two registered hypotheses, H3
and H6, concern quantized models. They are reported in a separate paper,
together with the conversion and on-device measurements. One logged deviation applies: the H5 source set and test details
were fixed after the freeze but before any EnCodec audio existed. The H2
resampling details, which the registration leaves open, were fixed after the
descriptive ITW rates had been seen. The cost-weighted check is registered, and
this paper is its first computation.

% ============================ 4. RESULTS ======================================
\section{Results}
\label{sec:results}

\begin{figure*}[!t]
\centering
\includegraphics[width=\textwidth]{fig_sensitivity.pdf}
\caption{Threshold sensitivity. HTER against the threshold, in standard
deviations of the target scores from the target-oracle threshold (vertical
line). Triangles mark the frozen \(\tauE\); the vertical distance from a
triangle to its curve's minimum is the TTG. A curve that is steep near its
minimum makes a small miss costly. fp32; LA19 dev thresholds.}
\label{fig:sens}
% src: analysis/figures/conference_figures.py (curves = paper_figures.collect(), checked equal
% to paper/figures/figures_data.json); plotted points in figures/conference_figures_data.json.
\end{figure*}

\subsection{Where the error lands}
% src: results/analysis/operating_points.txt (all FRR, FAR, TTG, EER below); safe bands from
% paper/figures/figures_data.json (= Table II "Band").
Table~\ref{tab:op} and Fig.~\ref{fig:frrfar}(a) give the operating behaviour on
every target. On LA19 eval the genuine speech comes from the source corpus and
only the attacks are new, so the error that grows is false acceptance. At
\(\tauE\) the ranked models accept 6.1--11.5\% of spoofs and reject only
0.1--0.5\% of genuine clips. The FRR budgets hold: at \(\tauF{1}\)
no model rejects more than 0.71\% of bona fide clips. Even the seen corpus is not
a safe harbour, though. AASIST and AASIST-L lose more than four HTER points to
unseen attacks (TTG 4.18 and 4.64\,pp, 95\% CIs [4.03, 4.33] and [4.49, 4.81]).
A threshold set on the evaluation set itself, as in \cite{zhou26_eer_hides},
would hide this part of the gap.

On ITW and DF21 the genuine speech has shifted, and the error moves to it. At
\(\tauE\), FRR exceeds FAR for four of five models on each target; the
exceptions are RawNet2 on ITW and AASIST on DF21. The FRR budgets do not
survive. A threshold that rejected 1\% of genuine dev speech rejects
42.8--81.8\% of genuine ITW speech and 18.1--55.0\% on DF21. At the 5\% budget,
ITW rejection reaches 75.5--98.2\%. Stricter budgets also widen the gap. Moving
from \(\tauE\) to \(\tauF{5}\) raises TTG for all five models on DF21 and for
all but RawNet2 on ITW, up to 37.0\,pp (XLS-R+AASIST on ITW).
% src (TTG at frr5 vs eer): results/analysis/h1_h4.json cells (ITW: 9.26/6.68/0.89/5.80/37.03
% vs 2.94/1.55/5.44/0.73/14.96; DF21: 4.25/6.68/2.15/6.90/18.90 vs 0.29/0.03/0.58/1.82/5.40).
The EER hides all of this. XLS-R+AASIST
reads 1.98\% EER on DF21 yet rejects 14.6\% of its genuine speech at \(\tauE\),
and reads 10.46\% on ITW yet rejects 49.9\%. On the clean-only DF21 surface
(305,525 trials), every DF21 statement here holds: FRR exceeds FAR for the same
four models, the 1\% budget becomes 14.0--49.8\%, and XLS-R+AASIST again has
the largest TTG at all three policies.
% src (clean-only): results/analysis/operating_points.txt, "== df21_clean" block and its
% "Stability" lines.

\begin{table}[!t]
\caption{Operating behaviour at the frozen LA19-dev thresholds (fp32, \%).
FRR, FAR and TTG at \(\tauE\); FRR\(_1\), FRR\(_5\): FRR at \(\tauF{1}\) and
\(\tauF{5}\), whose dev FRR is at most 1\% and 5\%; Band: 1-pp safe band (SD).}
\label{tab:op}
\centering
\footnotesize
\setlength{\tabcolsep}{3.2pt}
\begin{tabular}{@{}lrrrrrrr@{}}
\toprule
Model & EER & FRR & FAR & TTG & FRR\(_1\) & FRR\(_5\) & Band \\
\midrule
\multicolumn{8}{@{}l}{\emph{LA19 eval: seen corpus, unseen attacks}} \\
AASIST       &  0.83 &  0.16 &  9.83 &  4.18 &  0.29 &  3.58 & 0.68 \\
AASIST-L     &  0.99 &  0.10 & 11.14 &  4.64 &  0.23 &  4.51 & 0.62 \\
RawNet2      &  4.60 &  0.50 & 11.48 &  1.74 &  0.71 &  3.71 & 1.02 \\
Res-TSSDNet  &  1.62 &  0.26 &  6.08 &  1.58 &  0.34 &  2.05 & 0.92 \\
XLS-R (ref.) &  0.15 &  0.05 &  0.49 &  0.13 &  0.63 &  5.90 & 2.99 \\
\midrule
\multicolumn{8}{@{}l}{\emph{In-the-Wild}} \\
AASIST       & 43.02 & 43.77 & 41.91 &  2.94 & 57.42 & 98.23 & 0.62 \\
AASIST-L     & 44.46 & 50.65 & 36.54 &  1.55 & 64.43 & 97.34 & 0.82 \\
RawNet2      & 49.00 & 33.75 & 69.07 &  5.44 & 42.78 & 83.49 & 0.23 \\
Res-TSSDNet  & 34.17 & 50.28 & 18.97 &  0.73 & 54.37 & 75.52 & 0.89 \\
XLS-R (ref.) & 10.46 & 49.90 &  0.43 & 14.96 & 81.78 & 94.43 & 0.03 \\
\midrule
\multicolumn{8}{@{}l}{\emph{DF21, verified surface}} \\
AASIST       & 14.97 & 13.85 & 16.24 &  0.29 & 18.09 & 36.26 & 1.12 \\
AASIST-L     & 14.94 & 16.47 & 13.30 &  0.03 & 23.44 & 42.18 & 0.56 \\
RawNet2      & 19.65 & 33.68 &  5.60 &  0.58 & 35.17 & 41.00 & 3.57 \\
Res-TSSDNet  & 28.09 & 52.83 &  5.17 &  1.82 & 54.97 & 66.30 & 1.04 \\
XLS-R (ref.) &  1.98 & 14.56 &  0.16 &  5.40 & 29.11 & 41.72 & 0.20 \\
\bottomrule
\end{tabular}
% src: EER, FRR, FAR, TTG, FRR_1, FRR_5 = results/analysis/operating_points.txt (policy rows
% eer / frr1 / frr5); Band = safe_band_sd in paper/figures/figures_data.json. The TTG column
% equals results/analysis/h1_h4.json (ITW, DF21) to 1e-12 (checked by operating_points.py).
\end{table}

\begin{figure}[!t]
\centering
\includegraphics[width=\columnwidth]{fig_frr_far.pdf}
\caption{(a) FRR against FAR at the frozen \(\tauE\) for every model and target.
Points above the diagonal put the error on genuine speech. (b) The EnCodec
condition (H5) on the 1,000-clip manifest: each arrow runs from the clean copy
(filled) to its resynthesis (open).}
\label{fig:frrfar}
% src: figures/conference_figures_data.json (from results/analysis/operating_points.json
% and results/analysis/h5_test.json).
\end{figure}

\subsection{H2: false rejection dominates on ITW (confirmatory)}
% src: results/analysis/h2_test.txt and h2_test.json.
H2 is supported. False rejection exceeds false acceptance for three of the four
ranked models (Table~\ref{tab:tests}): Res-TSSDNet by 31.3\,pp, AASIST-L by
14.1\,pp and AASIST by 1.9\,pp, all with \(p<0.001\) and rejected under Holm
correction. Two caveats apply. The majority rests on AASIST's small margin,
which holds under five other bootstrap seeds and under unstratified resampling
(\(p\le0.002\)). RawNet2 fails in the opposite direction (\(-35.3\)\,pp): the
released checkpoint does not separate ITW at all (EER 49.0\%, AUC 0.51), and its
source threshold accepts most clips of either class. Outside the registered test,
the asymmetry grows at the stricter policies: at \(\tauF{5}\), FRR exceeds FAR
for all five models on ITW, RawNet2 included (83.5\% against 10.2\%).
% src (frr5 ITW): results/analysis/operating_points.txt, "== itw" block, frr5 rows.
% src (AUC 0.510): results/phase_2_conversion_app/main_matrix/itw_fp32_descriptives.txt.

\begin{table}[!t]
\caption{Confirmatory tests on the four ranked models (\%). One-sided
class-stratified paired bootstrap (\(B=1000\)); every rejection has
\(p<0.001\); Holm within each family.}
\label{tab:tests}
\centering
\footnotesize
\setlength{\tabcolsep}{3pt}
\begin{tabular}{@{}lccrc@{}}
\toprule
\multicolumn{5}{@{}l}{\emph{H2: ITW at \(\tauE\)}\hfill supported, 3 of 4} \\
Model & FRR & FAR & \(\frr-\far\) [95\% CI] & Rej. \\
\midrule
AASIST      & 43.77 & 41.91 & \(+1.86\) [\(+0.74\), \(+3.04\)] & yes \\
AASIST-L    & 50.65 & 36.54 & \(+14.11\) [\(+12.95\), \(+15.24\)] & yes \\
RawNet2     & 33.75 & 69.07 & \(-35.31\) [\(-36.38\), \(-34.28\)] & no \\
Res-TSSDNet & 50.28 & 18.97 & \(+31.32\) [\(+30.29\), \(+32.37\)] & yes \\
\midrule
\multicolumn{5}{@{}l}{\emph{H5: clean \(\to\) EnCodec 6\,kbps at \(\tauE\)}\hfill supported, 4 of 4} \\
Model & FRR & FAR & \(\Delta\frr-\Delta\far\) [95\% CI] & Rej. \\
\midrule
AASIST      & 0.2\(\to\)6.4  & 10.2\(\to\)3.8 & \(+12.6\) [\(+8.8\), \(+16.2\)] & yes \\
AASIST-L    & 0.0\(\to\)11.2 & 12.2\(\to\)5.8 & \(+17.6\) [\(+13.8\), \(+21.8\)] & yes \\
RawNet2     & 0.6\(\to\)1.8  & 14.2\(\to\)12.6 & \(+2.8\) [\(+1.0\), \(+4.6\)] & yes \\
Res-TSSDNet & 0.4\(\to\)68.4 & 5.0\(\to\)1.4 & \(+71.6\) [\(+67.0\), \(+75.8\)] & yes \\
\bottomrule
\end{tabular}
% src: results/analysis/h2_test.txt (top block); results/analysis/h5_test.txt (bottom block).
\end{table}

\subsection{Why some detectors lose more (H1, H4: exploratory)}
% src: results/analysis/h1_h4.txt (rho, CIs, TTG CIs); paper/figures/figures_data.json
% (tau_star_offset_sd, safe_band_sd); results/analysis/operating_points.txt (score percentiles;
% largest-TTG lines).
We find no detectable association between source accuracy and the gap. Over the
twelve ITW points (four ranked models, three policies), the Spearman correlation
between LA19 EER and TTG is \(\rho=0.00\) [\(-0.62\), \(+0.58\)]; with DF21 added (24 points) it is
\(+0.10\) [\(-0.32\), \(+0.53\)]. Only four distinct EERs underlie these points, so
they cluster by model. For H4, \(\rho(\log_{10}\text{params},\ttg)\) is \(+0.40\)
on ITW and \(+0.80\) on DF21, but with \(n=4\) both intervals span \([-1,1]\).
The 315M-parameter SSL reference, which stays outside every \(\rho\), has the best
EER on every target and the largest TTG on both shifted targets at all three
policies: 14.96\,pp [14.58, 15.39] on ITW and 5.40\,pp [5.10, 5.68] on DF21 at
\(\tauE\).

Fig.~\ref{fig:sens} shows why. The frozen threshold misses the target optimum
by between \(-1.64\)\,SD (RawNet2 on ITW) and \(+1.50\)\,SD (XLS-R+AASIST on
DF21), and the cost of a miss depends on the safe band. AASIST and AASIST-L
keep bands of 0.56--1.12\,SD on the shifted targets. The SSL reference's bands
are 0.20\,SD on DF21 and 0.03\,SD on ITW, so its moderate ITW miss
(\(+0.63\)\,SD) lands on a cliff. The cliff comes from saturation. On LA19 eval,
98\% of its spoof scores lie in \([-3.72,-2.80]\) and 98\% of its bona fide
scores in \([4.00,4.70]\) (1st--99th percentiles). \(\tauE=-1.75\) sits in the
empty gap between them, where any threshold is nearly perfect (band
2.99\,SD). On ITW the spoofs stay piled at the floor, with half of them in
\([-3.67,-3.61]\). The genuine speech, however, spreads over the whole range,
and its median (\(-1.74\)) falls on \(\tauE\); hence the 49.9\% FRR. The
target optimum (\(-3.54\)) sits on the upper edge of the spoof pile, and a
threshold only 0.13 units lower would accept three quarters of all spoofs. That
optimum lies inside the LA19 eval spoof range, where it would accept at least
10\% of those spoofs, so no threshold calibrated on source data would be placed
there.

Two cautions apply to reading TTG. A small gap is not good news when the ranking
itself has failed. On ITW the ranked models already sit at oracle HTERs of
34--46\%, so their curves have little valley to miss. A gap is also measured
against an oracle that the deployer never sees: it is the most that
re-calibration with target labels could recover.
% src: oracle HTER on ITW 39.89 / 42.04 / 45.97 / 33.89 = hter_oracle in
% results/analysis/operating_points.json (= figures_data.json).

\subsection{H5: a neural codec moves both classes (confirmatory)}
% src: results/analysis/h5_test.txt; DEVIATIONS.md DEV-2 (0.2-pp steps).
H5 is supported for all four ranked models (Table~\ref{tab:tests}). EnCodec
raises FRR more than FAR in every case, from \(+2.8\)\,pp (RawNet2) to
\(+71.6\)\,pp (Res-TSSDNet). The mechanism is broader than the hypothesis
states. False acceptance \emph{falls} for every model (Fig.~\ref{fig:frrfar}b):
resynthesis lowers the scores of genuine and fake speech alike. At a fixed
threshold that surfaces as rejected genuine speech, and Res-TSSDNet rejects
68.4\% of resynthesised bona fide clips. The test is met partly because
\(\Delta\far<0\), so both components should be reported. With 500 clips per
class, the rates move in 0.2-pp steps.

\subsection{Cost-weighted check (registered)}
% src: results/analysis/operating_points.txt (C columns; "Stability" lines) and
% operating_points.json (c_norm_min: ITW screening 1.000 / 1.000 / 1.000 / 0.994 / 0.861).
Under the screening model (\(\pi=0.95\)), every model at its transferred
threshold costs more than no detection on both shifted targets, at all three
policies. At \(\tauE\), \(C\) ranges from 7.1 to 10.0 on ITW and from 2.8 to
10.1 on DF21. On ITW the target-optimal threshold is no better: for AASIST,
AASIST-L and RawNet2 the cheapest choice is to accept everything (\(\min C=1.00\)).
Under the forensic model, XLS-R+AASIST is the only model below 1 on ITW at
\(\tauE\) (\(C=0.54\)), because its errors are almost all false rejections. The
ordering of the gaps is therefore not stable. XLS-R+AASIST has the largest HTER
gap on both shifted targets at all three policies. Under screening, the largest
cost gap belongs to AASIST-L, XLS-R+AASIST or AASIST on ITW (depending on the
policy) and to Res-TSSDNet on DF21. Under the forensic model it belongs to
RawNet2 on ITW and to AASIST or XLS-R+AASIST on DF21. The registration requires
an unstable conclusion to be reported as a finding, and this one is: the size of
a transfer gap is a property of the model \emph{and} the deployment, not of the
model alone.

% ============================ 5. DISCUSSION ==================================
\section{Discussion and Limitations}
\label{sec:discussion}
\emph{What to report.} The two failures found here, attacks accepted in-domain
and genuine speech rejected under shift, are invisible to EER. The second lands
on real people: a screening system with this behaviour flags genuine callers as
fake. We suggest that CM papers report, alongside EER:
\begin{enumerate}
\item the threshold the CM ships with, and the data and policy that set it;
\item FRR and FAR at that threshold on every target, at each operating point;
\item the safe-band width, which shows how much a small miss would cost;
\item the cost model under which any ranking of models is claimed.
\end{enumerate}
FRR budgets are fragile for a structural reason. They are set in the lower tail
of the source genuine-score distribution, and any downward shift of genuine
target scores moves mass across that tail first. Calibrating scores on source
data does not escape this. A calibrated score with a Bayes threshold
\cite{brummer06_calibration} is a monotone map followed by a source-chosen
threshold: one more point on the same curve. The results also suggest a
cheap check for deployers. False rejection, the dominant failure under shift,
can be estimated from genuine speech of the deployment channel alone, without
any attack data. Saturated detectors deserve particular care. A wide valley on
source data says nothing about where genuine target speech will fall, and the
unlabelled corrections that could move a threshold have proved unstable across
targets \cite{zhou26_eer_hides}.

\emph{Limitations.} (a) The ranked roster has four raw-waveform models from
2021--22 and a single graph-attention lineage. The one SSL model is an unranked
reference, so H1 and H4 rest on four models and their intervals are wide by
construction. (b) Every checkpoint is used as released, with one training run
each. The released RawNet2 does not separate ITW. (c) ITW has been the
community's transfer target since 2022 and is not an independent proxy for
deployment. Our DF21 figures use a verified surface and are not comparable with
published ones. (d) H5 covers one neural codec at one bitrate on 1,000 clips.
Conventional codecs, noise and reverberation are not tested here, and neither is
a language shift (e.g.\ MLAAD \cite{mlaad24}). (e) Everything here is fp32
desktop scoring.

% ============================ 6. CONCLUSION ==================================
\section{Conclusion}
\label{sec:conclusion}
A countermeasure calibrated once and deployed elsewhere fails asymmetrically.
Unseen attacks are accepted, but under corpus shift the larger failure is genuine
speech rejected, and false-rejection budgets set on source data do not survive. Source EER showed no
detectable association with the size of the failure. Under equal error costs,
the most accurate detector failed worst, because its saturated scores leave no
room for a threshold chosen in advance. A neural codec pushes genuine and fake
speech alike toward ``spoof''. These are properties of the
operating point, so they should be measured there: at a declared threshold, as
FRR and FAR, under a stated cost model. The harness, frozen thresholds, score
files and analysis scripts are released at \repo.

\bibliographystyle{IEEEtran}
\bibliography{refs}

\end{document}
